Calcul arithmétique sur un exemple : mille échanges possibles dont le gain s'étage régulièrement de 100 \euro{} à 0 \euro{}.\par Note : une taxe de t \euro{} fait disparaître les échanges qui rapportent moins de t \euro{}, soit 10 t échanges qui rapportaient en moyenne t/2 \euro{} ; la perte vaut 5 t². Les recettes valent t $\times$ (1 000 $-$ 10 t) : 9 000 \euro{} pour une taxe de 10 \euro{}, 16 000 \euro{} pour 20 \euro{}, 21 000 \euro{} pour 30 \euro{}. La perte, elle, passe de 500 à 2 000 puis 4 500 \euro{}.
